\documentclass[11pt,a4paper]{article}
\usepackage[ngerman]{babel}
\usepackage[TS1,T1]{fontenc}
\usepackage[utf8x]{inputenc}
\selectlanguage{ngerman}
\usepackage{theorem}
\usepackage[scaled=0.9]{helvet}
\usepackage{amsmath}
\usepackage{amssymb}
\usepackage{hyperref}
\usepackage{stmaryrd}
\usepackage{pgf,tikz}
\usepackage{relsize}
\usepackage{enumitem}
\usepackage{graphicx}
\usepackage{algpseudocode,amsmath,xifthen}

\newcounter{numb}
\theoremstyle{break}
\theorembodyfont{\upshape}
   	\newtheorem{exercise}{Exercise}[numb]
\setlength{\oddsidemargin}{0cm}
\setlength{\textwidth}{16cm}
\setlength{\textheight}{23cm}
\setlength{\topmargin}{-2cm}

\usetikzlibrary{shapes,arrows,automata,positioning,decorations.fractals}
\renewcommand\familydefault{\sfdefault}

\newcommand{\header}[2]{
\begin{minipage}[b]{\textwidth}
		\parbox[t]{5cm}{%
			\includegraphics[width=4cm]{unilogo}
			\hfill
		}
		\parbox[b]{11cm}{%
			%\scshape%
			Heinrich-Heine-University D\"usseldorf\\
			Computer Science Department\\
			Software Engineering and Programming Languages\\
			%Professor Dr.\ M.\ Leuschel
		Philipp K\"orner
}
		
\end{minipage}
	\begin{center}
		\bf
		Functional Programming -- WT 2024 / 2025\\
        Reading Guide {\thenumb}: {#1}
	\end{center}

    \pagestyle{empty}
    \paragraph{Timeline:} This unit should be completed by {#2}.
}

\setcounter{numb}{9}


\begin{document}
	
\begin{minipage}[b]{\textwidth}
	\parbox[t]{5cm}{%
		\includegraphics[width=4cm]{unilogo}
		\hfill
	}
	\parbox[b]{11cm}{%
		%\scshape%
		Heinrich-Heine-University D\"usseldorf\\
		Computer Science Department\\
		Software Engineering and Programming Languages\\
		%Professor Dr.\ M.\ Leuschel
		Philipp K\"orner
	}
\end{minipage}
\begin{center}
	\bf
	Functional Programming -- WT 2023 / 2024\\
	Reading Guide 9: Time and Value (Theory)
\end{center}

\pagestyle{empty}

\paragraph{Timeline:} This unit should be completed by 11.12.2023.

\section{Material} 

\begin{itemize}
\item Rich Hickey: Are We There Yet? \\ \url{https://www.infoq.com/presentations/Are-We-There-Yet-Rich-Hickey/}
\item Rich Hickey: Clojure Concurrency \url{https://youtu.be/nDAfZK8m5_8} (21:06 -- 1:16:05 --- first twenty minutes may be used as a revision)
\item Learning Video: Epochal Time Model: \url{https://mediathek.hhu.de/watch/20e935c4-9934-4d4d-abe7-9bc3925ff67d}
\end{itemize}

\paragraph{Note:} This unit only deals with the theoretical foundations.
In two additional reading guides, we will inspect the semantics of different implementations of this model
in more detail (and code).

\section{Learning Outcomes}

After completing this unit you should be able to

\begin{itemize}
	\item define and assign the terms identity, value observer and state.
    \item describe the epochal time-model and apply it to different implementations of state management.
    \item describe compare-and-swap (CAS)-semantics\footnote{You may refer to a Wikipedia page on CAS.}.
    \item describe the role of persistent data structure in a concurrent context.
\end{itemize}


\section{Exercises}

No practical exercises for a theoretical introduction.
Proceed with the usual 4clojure exercises.


%\begin{exercise}[Higher Order Functions]
%\begin{enumerate}[label=\alph*)]
%
% \item Write a function \texttt{flip}, which reverses the arguments of a function.
% Examples:
% \begin{verbatim}
%user=> (nth [3 4 5 6] 2)
%5
%user=> ((flip nth) 2 [3 4 5 6])
%5
%user=> (- 3 2 1)
%0
%user=> ((flip -) 1 2 3)
%0
%user=> (drop 2 [4 5 6 7])
%(6 7)
%user=> ((flip drop) [4 5 6 7] 2)
%(6 7) \end{verbatim}
%\item The Clojure function \texttt{comp} can be used to generate a composition of functions. The result of the call \texttt{((comp f g) x)} is the same as the the result of the call \texttt{(f (g x)))}. Implement your own version of \texttt{comp}. Your version should only cover the most general case \texttt{(defn mycomp [\& fs] ...)} instead of optimizing for various smaller arities like the Clojure version. 
%
%\begin{verbatim}
%user=> ((mycomp inc (fn [x] (* x x))) 4)
%17
%user=> ((mycomp (fn [x] (* x x)) inc) 4)
%25 \end{verbatim}
%
%
%
%\item Implement the Clojure function \texttt{juxt}, which receives functions as parameters and creates a new function. This function receives an input, applies each function to the input and creates a vector with the results. Implement the most general case of \texttt{juxt}.
%
%\begin{verbatim}
%user=> ((myjuxt inc dec (fn [x] (* x x))) 3)
%[4 2 9] \end{verbatim}
%
%\end{enumerate}
%\end{exercise}
%
%
%
%\begin{exercise}[\texttt{every?}]
%    \texttt{(fn every? [f c])} is a higher-order function, which receives a function $f$ and a collection $c$ as parameters. 
%    The return value is \texttt{true}, if the function $f$ returns a truthey value for every element in the collection and otherwise \texttt{false}.
%
%    \begin{verbatim}
%(every? even? [1 2 3 4]) 
%=> false
%(every? even? [2 4])
%=> true \end{verbatim}
%
%    Implement \texttt{every?} yourself.
%\end{exercise}


	\section*{Questions}
	If you have any questions, please contact Philipp K"orner (\texttt{p.koerner@hhu.de}) or post it to the Rocket.Chat group.
\end{document}

